\documentclass[11pt,a4paper]{article}

\usepackage[a4paper,margin=2.5cm]{geometry}
\usepackage{amsmath}
\usepackage{amssymb}
\usepackage{graphicx}
\usepackage{hyperref}
\usepackage{microtype}
\raggedbottom

\title{Image Analysis Coursework Report}
\author{eid23}
\date{MPhil in Data Intensive Science\\Submission: 24 June 2026}

\begin{document}
\maketitle

\section{Part I: Inverse Problems in Imaging}

\subsection{A.1 Forward Problem: Simulating Image Acquisition}

The first stage models the acquisition process for a diffraction-limited optical
instrument. Starting from a clean microscopy image \(g \in \mathbb{R}^{m \times n}\),
the observed image is generated as
\begin{equation}
    f = H g + w,
\end{equation}
where \(H\) is a point-spread-function blur operator and \(w\) is additive
Gaussian noise. The implementation normalises the selected TIFF frame to
floating point intensities in \([0,1]\), applies a Gaussian PSF with standard
deviation \(\sigma = 2.0\), and then adds zero-mean Gaussian noise with
\(\sigma_n = 0.001\). The random number generator is seeded with 42 by default,
so the simulated acquisition is reproducible.

The default command downloads the PhC-C2DH-U373 microscopy dataset, selects the
first sorted frame from sequence 01 as the ground truth image, and writes the
ground truth, blurred image, blurred noisy image, comparison figure, and metadata
to \texttt{outputs/part\_i/a1/}.

\begin{figure}[htbp]
    \centering
    \IfFileExists{../outputs/part_i/a1/comparison.png}{%
        \includegraphics[width=\linewidth]{../outputs/part_i/a1/comparison.png}
    }{%
        \fbox{\parbox{0.9\linewidth}{%
            Part I.A.1 figure generated by \texttt{python scripts/run\_a1.py}.
        }}
    }
    \caption{Forward acquisition simulation: clean ground truth, Gaussian-blurred
    image, and blurred image with additive Gaussian noise.}
    \label{fig:a1-forward}
\end{figure}

\subsection{A.2 Classical Image Restoration}

The second stage treats the simulated observation from Part I.A.1 as the input
for classical restoration. I implemented isotropic total variation denoising by
approximately solving
\begin{equation}
    \min_g \frac{1}{2}\lVert g - f \rVert_2^2 + \lambda\,\mathrm{TV}(g),
\end{equation}
where the discrete TV term is computed from forward finite differences in the
horizontal and vertical directions. The optimisation is solved with a
Chambolle--Pock primal-dual method~\cite{chambollepock}. The dual variables are projected onto the
pointwise \(\ell_2\) ball of radius \(\lambda\), and the objective value is
recorded after each iteration.

I evaluate three regularisation strengths by default:
\(\lambda=0.001\), \(\lambda=0.01\), and \(\lambda=0.08\). These are intended to
represent under-regularisation, a moderate compromise, and over-regularisation.
Small \(\lambda\) values keep the solution close to the measured image, so fine
texture and noise are both retained. Larger \(\lambda\) values suppress noise
more aggressively but can flatten subtle structures. On the default frame, the
noisy blurred observation gives PSNR 33.91 dB and SSIM 0.9219. A Gaussian filter
with \(\sigma=1\) gives PSNR 33.35 dB and SSIM 0.9154, while TV denoising gives
PSNR/SSIM values of 33.86/0.9201, 33.50/0.9094, and 31.90/0.8848 for
\(\lambda=0.001\), \(0.01\), and \(0.08\), respectively. These values show that
the strongest regularisation is over-smoothing this example. TV regularisation
tends to preserve edges better than a standard Gaussian filter because the
penalty is on the magnitude of the image gradient rather than on local averaging
itself; sharp transitions can remain when they are supported by the data
fidelity term. Here the input noise is already very small, so every additional
smoothing operation reduces PSNR. Moreover, the stated A.2 objective contains
no forward operator \(H\): it is a denoising model and cannot invert the
Gaussian blur. Thus \(\lambda=0.001\) is best numerically, while \(0.01\)
illustrates a stronger visual noise/detail compromise rather than genuine
deconvolution.

\begin{figure}[htbp]
    \centering
    \IfFileExists{../outputs/part_i/a2/comparison.png}{%
        \includegraphics[width=\linewidth]{../outputs/part_i/a2/comparison.png}
    }{%
        \fbox{\parbox{0.9\linewidth}{%
            Part I.A.2 figure generated by \texttt{python scripts/run\_a2.py}.
        }}
    }
    \caption{Classical restoration comparison for the noisy blurred input,
    Gaussian filtering, and TV denoising at several regularisation strengths.}
    \label{fig:a2-restoration}
\end{figure}

\begin{figure}[htbp]
    \centering
    \IfFileExists{../outputs/part_i/a2/objective.png}{%
        \includegraphics[width=0.75\linewidth]{../outputs/part_i/a2/objective.png}
    }{%
        \fbox{\parbox{0.75\linewidth}{%
            Part I.A.2 objective trace generated by \texttt{python scripts/run\_a2.py}.
        }}
    }
    \caption{Objective value during the TV denoising iterations. A decreasing
    trace indicates that the iterative solver is moving toward a minimiser of
    the regularised objective.}
    \label{fig:a2-objective}
\end{figure}

\subsection{B(i) Algorithm Unrolling}

For the data-driven reconstruction experiment, I implemented a learned unrolled
proximal-gradient network. The forward operator \(H\) is a Gaussian blur with an
\(11\times 11\) kernel and standard deviation \(2.0\). Training observations are
simulated as \(f = Hg + w\), with white Gaussian noise of variance \(0.001\). The
network uses \(K=6\) unrolled iterations by default:
\begin{equation}
    g^{(k+1)} = \phi_{\theta^{(k)}}\left(g^{(k)} - \eta^{(k)}H^\top(Hg^{(k)} - f)\right),
\end{equation}
where each \(\phi_{\theta^{(k)}}\) is a small CNN proximal map with two hidden
convolutional layers and PReLU activations. The step sizes are learned as
unconstrained parameters passed through a softplus transform, which keeps
\(\eta^{(k)}\geq 0\) during training. The training loss is the \(\ell_1\) error
between the final reconstruction \(g^{(K)}\) and the clean reference image.
Training uses random crops and random horizontal/vertical flips from sequence 01;
sequence 02 is used for test metrics.

The full GPU run used 50 epochs, 32 feature channels, batch size 4,
\(128\times128\) crops and learning rate \(10^{-4}\). On all 115 images in
sequence 02, the degraded input scored PSNR 28.18 dB and SSIM 0.4897, while the
unrolled reconstruction improved this to PSNR 34.10 dB and SSIM 0.9113.

\begin{figure}[htbp]
    \centering
    \IfFileExists{../outputs/part_i/b1_full/comparison.png}{%
        \includegraphics[width=\linewidth]{../outputs/part_i/b1_full/comparison.png}
    }{%
        \fbox{\parbox{0.9\linewidth}{%
            Part I.B(i) comparison generated by \texttt{python scripts/run\_b1.py}.
        }}
    }
    \caption{Visual comparison for the unrolled PGD model: ground truth, blurred
    noisy input, and learned reconstruction.}
    \label{fig:b1-unrolling}
\end{figure}


\subsection{B(ii) Out-of-Distribution Generalisation}

For the OOD experiment the same independent-layer unrolled PGD architecture is
trained with a more severe degradation, using a Gaussian blur kernel of size
\(21\times21\), standard deviation \(4.0\), and white Gaussian noise of variance
\(0.001\). At test time the model is evaluated on the original Part I.B(i)
setting, namely an \(11\times11\) Gaussian kernel with standard deviation \(2.0\)
and the same noise variance. This isolates whether the learned proximal maps are
robust to a forward model mismatch. The script records average PSNR and SSIM for
both the degraded input and the reconstruction, plus a visual comparison. This
tests whether training on stronger blur makes the learned prior robust or makes
it overcompensate when the test blur is weaker.

The full OOD run achieved PSNR 33.11 dB and SSIM 0.9030 on the B(i) test
degradation. Relative to the matched B(i) model, this is a reduction of 0.99 dB
and 0.0083 SSIM. The learned solver is therefore sensitive to a mismatch between
the training and test forward operators, despite training on the more severe blur.

\begin{figure}[htbp]
    \centering
    \IfFileExists{../outputs/part_i/b2_full/comparison.png}{%
        \includegraphics[width=\linewidth]{../outputs/part_i/b2_full/comparison.png}
    }{%
        \fbox{\parbox{0.9\linewidth}{Part I.B(ii) comparison generated by \texttt{python scripts/run\_b2.py}.}}
    }
    \caption{Part I.B(ii) OOD unrolled PGD comparison.}
    \label{fig:b2-ood}
\end{figure}

\subsection{B(iii) Shared-Parameter Convergence}

For the convergence experiment I train a shared-parameter unrolled PGD model:
\begin{equation}
    g^{(k+1)} = \psi_{\theta}\left(g^{(k)} - \eta H^\top(Hg^{(k)} - f)\right),
\end{equation}
with a single CNN proximal map and a single positive softplus-constrained step
size reused at every layer. Training uses \(K=6\) iterations under the matched
\(11\times11\), \(\sigma=2.0\) blur and noise variance \(0.001\). At evaluation,
the same learned update is applied to one test image for \(T=K,4K,8K,16K\). The
script writes a PSNR/SSIM table and a PSNR curve. The measured values are:

\begin{center}
\begin{tabular}{c|rrrr}
\(T\) & 6 & 24 & 48 & 96 \\ \hline
PSNR (dB) & 33.61 & 23.88 & 18.05 & 6.56 \\
SSIM & 0.9049 & 0.8278 & 0.5773 & 0.0068
\end{tabular}
\end{center}
The monotonic deterioration beyond \(K\) shows that the finite-depth training
has not learned a globally stable fixed point. The shared update is effective
at its trained depth but should not be interpreted as a convergent solver.

\begin{figure}[htbp]
    \centering
    \IfFileExists{../outputs/part_i/b3_full/convergence_psnr.png}{%
        \includegraphics[width=0.7\linewidth]{../outputs/part_i/b3_full/convergence_psnr.png}
    }{%
        \fbox{\parbox{0.75\linewidth}{Part I.B(iii) convergence plot generated by \texttt{python scripts/run\_b3.py}.}}
    }
    \IfFileExists{../outputs/part_i/b3_full/comparison.png}{%
        \includegraphics[width=\linewidth]{../outputs/part_i/b3_full/comparison.png}
    }{}
    \caption{Part I.B(iii) PSNR and reconstructions as the shared update is
    iterated beyond the training depth.}
    \label{fig:b3-convergence}
\end{figure}

\section{Part II: Motion Estimation and Cell Tracking}

\subsection{A. Motion Estimation}

\paragraph{(i) Pair and change visualisation.}
I use consecutive sequence-01 frames \texttt{t000} and \texttt{t001}, acquired
15 minutes apart. Both are normalised to \([0,1]\). Their absolute difference
highlights changes at cell boundaries and within textured cell interiors while
leaving most of the substrate dark (Figure~\ref{fig:motion-summary}).

\paragraph{(ii--iii) Horn--Schunck implementation and test.}
Spatial derivatives are computed from the mean of the two frames and
\(D_t f=f_{t+1}-f_t\). Following Horn and Schunck~\cite{hornschunck}, the
implemented Jacobi iteration uses four-neighbour averages and minimises
\begin{equation}
 E(u,v)=\langle(D_xfu+D_yfv+D_tf)^2\rangle+
 \alpha^2\langle|\nabla u|^2+|\nabla v|^2\rangle .
\end{equation}
Tests translate a synthetic square by one pixel in each cardinal direction.
The estimated component has the correct sign and the orthogonal component is
near zero; its magnitude is attenuated by smoothness and the aperture problem.

\paragraph{(iv) Energy.}
For \(\alpha=1\) and 100 iterations, energy falls from
\(1.8590\times10^{-3}\) to a minimum \(1.8395\times10^{-3}\) at iteration 49,
then rises slightly to \(1.8416\times10^{-3}\). The Jacobi sequence is therefore
not monotonically energy decreasing; iteration 49 is a better stopping point
than blindly using the last iterate.

\paragraph{(v) Estimated flow.}
HSV hue encodes direction and brightness encodes magnitude. Motion is
concentrated around phase-contrast edges and intracellular texture, while the
background is nearly stationary. The background-inclusive mean vector is
\((\bar u,\bar v)=(-0.0026,-0.0043)\) pixels per frame and mean magnitude is
0.0197 pixels; arrows in Figure~\ref{fig:motion-summary} are enlarged 100 times.

\begin{figure}[htbp]
 \centering
 \IfFileExists{../outputs/part_ii/a/motion_summary.png}{%
  \includegraphics[width=\linewidth]{../outputs/part_ii/a/motion_summary.png}
 }{\fbox{\parbox{0.9\linewidth}{Generated by \texttt{python scripts/run\_motion.py}.}}}
 \caption{Input pair, absolute difference, HSV and quiver flow, and energy.
 The dashed line marks the minimum-energy iterate.}
 \label{fig:motion-summary}
\end{figure}

\paragraph{(vi) Regularisation and iterations.}
The sweep in Figure~\ref{fig:motion-sweep} compares \(\alpha=0.3\), 1 and 3.
Iteration counts are 10, 50, 100 and 200. At 100 iterations,
mean magnitude falls from 0.110 pixels for \(\alpha=0.3\), to 0.0197 for
\(\alpha=1\), and 0.00257 for \(\alpha=3\): stronger regularisation suppresses
and smooths motion. Additional iterations initially propagate motion into
neighbouring pixels, but all settings eventually pass their measured energy
minimum. This motivates parameter validation and energy-based stopping.

\begin{figure}[htbp]
 \centering
 \includegraphics[width=0.82\linewidth]{../outputs/part_ii/a/parameter_sweep.png}
 \caption{Sensitivity of final energy and mean flow magnitude to \(\alpha\) and
 iteration count.}
 \label{fig:motion-sweep}
\end{figure}

\subsection{B. YOLO Cell Segmentation}

\paragraph{(i) Export.}
All 34 annotated frames are exported: 15 sequence-01 frames for training and
19 sequence-02 frames for validation. Each labelled component is converted to
its longest simplified contour and coordinates are normalised by image width
and height. Rasterising the polygons back to binary masks gives mean IoU 0.9967
and minimum IoU 0.9769. The model must distinguish whole cells from
phase-contrast halos, intracellular texture, debris, substrate features, and
ambiguous partial objects at image boundaries.

\paragraph{(ii) Training.}
A pretrained \texttt{yolo11n-seg.pt} model~\cite{ultralytics} was trained on a
GPU for 100 epochs with \(640\times640\) inputs and batch size 16. On sequence
02, segmentation mAP50--95(M) is 0.694. This estimate is based on only 19
validation images and is not exhaustive temporal validation.

\paragraph{(iii) Label-image conversion.}
Inference polygons are filled with \texttt{draw.polygon}; each accepted polygon receives a distinct
positive integer and background remains zero. The confidence threshold is 0.25.

\paragraph{(iv--v) Frame choice and quality.}
I inspect unannotated early, middle and late frames \texttt{t000},
\texttt{t057}, and \texttt{t114}, testing temporal coverage without simply
replotting training annotations. Figure~\ref{fig:yolo} shows only 2, 3 and 2
masks respectively despite additional visible cell-like objects. Boundaries of
detected large cells are plausible, but recall is poor for smaller, brighter
and partial cells. The validation mAP therefore masks meaningful under-detection
on the complete sequence, likely due to sparse annotations.

\begin{figure}[htbp]
 \centering
 \IfFileExists{../outputs/part_ii/b/yolo_dataset/verify_overlay.png}{%
  \includegraphics[width=0.28\linewidth]{../outputs/part_ii/b/yolo_dataset/verify_overlay.png}
 }{}
 \IfFileExists{../outputs/part_ii/b/prediction_montage.png}{%
  \includegraphics[width=0.7\linewidth]{../outputs/part_ii/b/prediction_montage.png}
 }{}
 \caption{Export verification (left) and predictions on early, middle and late
 unannotated sequence-01 frames (right).}
 \label{fig:yolo}
\end{figure}

\subsection{C. Cell Tracking}

\paragraph{(i) Sequence segmentation.}
The trained model is loaded once and applied at confidence 0.25 to all 115
sequence-01 frames. It produces 287 masks and no mask in one frame.

\paragraph{(ii) Centroids and missing information.}
\texttt{regionprops} supplies each mask's centroid and area, but not temporal
identity: a tracker must infer which detection in frame \(t+1\) belongs to a
cell from frame \(t\), while handling missed, merged and split masks.

\paragraph{(iii) LapTrack parameters.}
The completed run uses \texttt{LapTrack}~\cite{laptrack}, with squared Euclidean
linking and gap-closing costs, a 35-pixel cutoff (\(35^2=1225\)), one
missing-frame gap, and splitting/merging disabled. At 0.65 micrometres per pixel
this permits 22.75 micrometres per 15-minute interval. It exceeds the observed
95th-percentile and maximum linked speeds (12.0 and 24.2 pixels/frame), providing
margin for fast cells, although a permissive cutoff can create incorrect links.
A one-frame gap can bridge one missing mask; longer omissions terminate a track.

\paragraph{(iv) Visualisation and limitations.}
LapTrack returns 43 tracks, with median length 3 frames, maximum length 29, and
23 tracks lasting at most three frames (Figure~\ref{fig:tracking}). This severe
fragmentation is consistent with YOLO under-detection and is more informative
than the total track count alone.

\begin{figure}[htbp]
 \centering
 \includegraphics[width=\linewidth]{../outputs/part_ii/c_yolo/tracking_summary.png}
 \caption{Full trajectories projected on \texttt{t000}, detections per frame,
 and track-length distribution. Circles and crosses mark starts and ends.}
 \label{fig:tracking}
\end{figure}

\paragraph{(v) Relationship to optical flow.}
For each linked cell and frame interval, I would average the dense flow inside
its mask and compare that internal-motion vector with the centroid displacement.
Their angle measures alignment and their magnitudes separate intracellular
reorganisation from whole-cell translation. Uncertain or missing masks should
be excluded or propagated with the track model.

\begin{thebibliography}{9}
\bibitem{chambollepock}
A. Chambolle and T. Pock, A first-order primal-dual algorithm for convex
problems with applications to imaging, \emph{J. Math. Imaging Vis.}, 2011.
\bibitem{hornschunck}
B. Horn and B. Schunck, Determining optical flow, \emph{Artificial
Intelligence}, 1981.
\bibitem{ultralytics}
Ultralytics, \emph{YOLO segmentation documentation},
\url{https://docs.ultralytics.com/tasks/segment/}.
\bibitem{laptrack}
LapTrack developers, \emph{LapTrack documentation},
\url{https://laptrack.readthedocs.io/}.
\end{thebibliography}

\end{document}
